% Options for packages loaded elsewhere
\PassOptionsToPackage{unicode}{hyperref}
\PassOptionsToPackage{hyphens}{url}
\documentclass[
  ignorenonframetext,
]{beamer}
\newif\ifbibliography
\usepackage{pgfpages}
\setbeamertemplate{caption}[numbered]
\setbeamertemplate{caption label separator}{: }
\setbeamercolor{caption name}{fg=normal text.fg}
\beamertemplatenavigationsymbolsempty
% remove section numbering
\setbeamertemplate{part page}{
  \centering
  \begin{beamercolorbox}[sep=16pt,center]{part title}
    \usebeamerfont{part title}\insertpart\par
  \end{beamercolorbox}
}
\setbeamertemplate{section page}{
  \centering
  \begin{beamercolorbox}[sep=12pt,center]{section title}
    \usebeamerfont{section title}\insertsection\par
  \end{beamercolorbox}
}
\setbeamertemplate{subsection page}{
  \centering
  \begin{beamercolorbox}[sep=8pt,center]{subsection title}
    \usebeamerfont{subsection title}\insertsubsection\par
  \end{beamercolorbox}
}
% Prevent slide breaks in the middle of a paragraph
\widowpenalties 1 10000
\raggedbottom
\AtBeginPart{
  \frame{\partpage}
}
\AtBeginSection{
  \ifbibliography
  \else
    \frame{\sectionpage}
  \fi
}
\AtBeginSubsection{
  \frame{\subsectionpage}
}
\usepackage{iftex}
\ifPDFTeX
  \usepackage[T1]{fontenc}
  \usepackage[utf8]{inputenc}
  \usepackage{textcomp} % provide euro and other symbols
\else % if luatex or xetex
  \usepackage{unicode-math} % this also loads fontspec
  \defaultfontfeatures{Scale=MatchLowercase}
  \defaultfontfeatures[\rmfamily]{Ligatures=TeX,Scale=1}
\fi
\usepackage{lmodern}
\ifPDFTeX\else
  % xetex/luatex font selection
\fi
% Use upquote if available, for straight quotes in verbatim environments
\IfFileExists{upquote.sty}{\usepackage{upquote}}{}
\IfFileExists{microtype.sty}{% use microtype if available
  \usepackage[]{microtype}
  \UseMicrotypeSet[protrusion]{basicmath} % disable protrusion for tt fonts
}{}
\makeatletter
\@ifundefined{KOMAClassName}{% if non-KOMA class
  \IfFileExists{parskip.sty}{%
    \usepackage{parskip}
  }{% else
    \setlength{\parindent}{0pt}
    \setlength{\parskip}{6pt plus 2pt minus 1pt}}
}{% if KOMA class
  \KOMAoptions{parskip=half}}
\makeatother
\usepackage{longtable,booktabs,array}
\newcounter{none} % for unnumbered tables
\usepackage{calc} % for calculating minipage widths
\usepackage{caption}
% Make caption package work with longtable
\makeatletter
\def\fnum@table{\tablename~\thetable}
\makeatother
\usepackage{graphicx}
\makeatletter
\newsavebox\pandoc@box
\newcommand*\pandocbounded[1]{% scales image to fit in text height/width
  \sbox\pandoc@box{#1}%
  \Gscale@div\@tempa{\textheight}{\dimexpr\ht\pandoc@box+\dp\pandoc@box\relax}%
  \Gscale@div\@tempb{\linewidth}{\wd\pandoc@box}%
  \ifdim\@tempb\p@<\@tempa\p@\let\@tempa\@tempb\fi% select the smaller of both
  \ifdim\@tempa\p@<\p@\scalebox{\@tempa}{\usebox\pandoc@box}%
  \else\usebox{\pandoc@box}%
  \fi%
}
% Set default figure placement to htbp
\def\fps@figure{htbp}
\makeatother
\setlength{\emergencystretch}{3em} % prevent overfull lines
\providecommand{\tightlist}{%
  \setlength{\itemsep}{0pt}\setlength{\parskip}{0pt}}
\usepackage{bookmark}
\IfFileExists{xurl.sty}{\usepackage{xurl}}{} % add URL line breaks if available
\urlstyle{same}
\hypersetup{
  hidelinks,
  pdfcreator={LaTeX via pandoc}}

\author{\texorpdfstring{}{}}
\date{}

\begin{document}

\begin{frame}
\newpage
\end{frame}

\section{Discussion}\label{sec:discussion}

\begin{frame}{Synthesis of Findings}
\protect\phantomsection\label{synthesis-of-findings}
Our simulation-based evaluation across six multiagent architecture
models validates the core theoretical claims of the Cognitive Integrity
Framework established in Part 1. The 94\% overall detection rate
achieved by the full CIF deployment represents a substantial improvement
over any individual defense mechanism, confirming that the
multiplicative composition theorems translate from formal proofs to
practical protection. More importantly, the consistency of this result
across architecturally diverse systems---from hierarchical orchestrator
patterns to peer-to-peer topologies---suggests that CIF's formal
abstractions capture genuine structural properties of multiagent
security rather than artifacts of specific implementation choices. We
now examine these findings in detail, beginning with the mechanisms
underlying layered defense success.

\begin{block}{Why Layered Defense Succeeds}
\protect\phantomsection\label{why-layered-defense-succeeds}
\begin{figure}
\centering
\includegraphics[width=0.95\linewidth,height=\textheight,keepaspectratio,alt={Defense Composition Architecture. Diagram illustrating the series and parallel composition of CIF defense mechanisms. The Cognitive Firewall provides the first line of defense (input filtering), followed by the Belief Sandbox (provisional isolation) and Tripwires (continuous monitoring) in series. Trust Calculus and Byzantine Consensus operate in parallel for delegation and coordination decisions. The multiplicative detection guarantee (Part 1, Theorems 3.1-3.2) emerges from the orthogonality of attack surfaces targeted by each layer: series composition yields P\_\{\textbackslash\textbackslash text\{detect\}\} = 1 - \textbackslash\textbackslash prod\_\{i\}(1 - r\_i) computed via \textbackslash texttt\{compute\textbackslash\_series\textbackslash\_detection\textbackslash\_rate()\}, while parallel composition uses max-score fusion via \textbackslash texttt\{compute\textbackslash\_parallel\textbackslash\_detection\textbackslash\_rate()\}. The Venn overlap statistics in the accompanying table are dynamically computed from per-mechanism detection rates using these composition functions.}]{../figures/defense_composition.pdf}
\caption{Defense Composition Architecture. Diagram illustrating the
series and parallel composition of CIF defense mechanisms. The Cognitive
Firewall provides the first line of defense (input filtering), followed
by the Belief Sandbox (provisional isolation) and Tripwires (continuous
monitoring) in series. Trust Calculus and Byzantine Consensus operate in
parallel for delegation and coordination decisions. The multiplicative
detection guarantee (Part 1, Theorems 3.1-3.2) emerges from the
orthogonality of attack surfaces targeted by each layer: series
composition yields \(P_{\\text{detect}} = 1 - \\prod_{i}(1 - r_i)\)
computed via
\textbackslash texttt\{compute\textbackslash\_series\textbackslash\_detection\textbackslash\_rate()\},
while parallel composition uses max-score fusion via
\textbackslash texttt\{compute\textbackslash\_parallel\textbackslash\_detection\textbackslash\_rate()\}.
The Venn overlap statistics in the accompanying table are dynamically
computed from per-mechanism detection rates using these composition
functions.}\label{fig:defense-composition}
\end{figure}

The defense composition architecture (\cref{fig:defense-composition})
illustrates how the CIF defense mechanisms integrate into a coherent
defense posture. The multiplicative composition of detection rates
(Theorems 3.1-3.2 in Part 1) explains the empirical observation that
full CIF substantially outperforms individual mechanisms. Each defense
targets a distinct attack surface:

{\def\LTcaptype{none} % do not increment counter
\begin{longtable}[]{@{}lll@{}}
\toprule\noalign{}
Defense Layer & Target Attack Surface & Contribution \\
\midrule\noalign{}
\endhead
Cognitive Firewall & Input-based injection & Blocks direct attacks \\
Belief Sandbox & Unverified content & Contains propagation \\
Tripwires & Belief manipulation & Detects subtle drift \\
Trust Calculus & Delegation abuse & Bounds amplification \\
Consensus & Coordination attacks & Ensures agreement integrity \\
\bottomrule\noalign{}
\end{longtable}
}
\end{block}

\begin{block}{Architecture-Specific Insights}
\protect\phantomsection\label{architecture-specific-insights}
\textbf{Table: Architecture vulnerability patterns and observed CIF
defense response.} \{\#tab:architecture-insights\}

{\def\LTcaptype{none} % do not increment counter
\begin{longtable}[]{@{}
  >{\raggedright\arraybackslash}p{(\linewidth - 4\tabcolsep) * \real{0.3333}}
  >{\raggedright\arraybackslash}p{(\linewidth - 4\tabcolsep) * \real{0.3333}}
  >{\raggedright\arraybackslash}p{(\linewidth - 4\tabcolsep) * \real{0.3333}}@{}}
\toprule\noalign{}
\begin{minipage}[b]{\linewidth}\raggedright
Architecture
\end{minipage} & \begin{minipage}[b]{\linewidth}\raggedright
Primary Vulnerability
\end{minipage} & \begin{minipage}[b]{\linewidth}\raggedright
Observed CIF Defense Mechanism
\end{minipage} \\
\midrule\noalign{}
\endhead
Hierarchical & Orchestrator compromise cascades & Orchestrator-specific
tripwires (82\% detection) \\
Peer-to-peer & Lateral movement amplification & Byzantine consensus +
trust decay \\
Role-based & Role impersonation & Attestation verification at role
transitions \\
State machine & State corruption & State hash verification
(deterministic detection) \\
\bottomrule\noalign{}
\end{longtable}
}

The architecture-specific results reveal that vulnerability patterns
align closely with the structural properties predicted by Part 1's
threat model analysis. Hierarchical architectures concentrate risk at
the orchestrator: a single compromised orchestrator can cascade
malicious instructions to all subordinate agents. Our evaluation shows
that tripwire-only deployments achieve 82\% detection in hierarchical
topologies but only 61\% in peer-to-peer systems, quantifying the
architectural dependence of defense effectiveness.

Peer-to-peer architectures present the opposite profile. Without a
central authority, lateral movement between agents is the primary threat
vector. Trust amplification through delegation chains enables an
attacker who compromises a single agent to gradually extend influence
across the network. The trust calculus with \(\delta^d\) decay directly
addresses this: the exponential decay bound ensures that delegated trust
diminishes with chain length, preventing unbounded amplification. Our
results confirm that peer-to-peer topologies show the largest relative
improvement (from 0\% baseline to 94\% with full CIF), consistent with
the theoretical prediction that these architectures benefit most from
formal trust bounds.

Role-based systems introduce impersonation as the primary risk. When
agents assume specialized roles (researcher, writer, reviewer), an
attacker who can assume a trusted role gains the permissions associated
with that role. In our evaluation, attestation-based verification at
role transitions detected 94\% of impersonation attempts
(\cref{tab:arch-ci}). Unexpected role transitions served as reliable
early indicators of compromise.
\end{block}
\end{frame}

\begin{frame}{Limitations and Threats to Validity}
\protect\phantomsection\label{limitations-and-threats-to-validity}
\begin{block}{Residual Attack-Type Vulnerabilities}
\protect\phantomsection\label{residual-attack-type-vulnerabilities}
Despite strong overall performance, specific attack types remain
challenging and merit detailed examination.

Semantic equivalent attacks pose the most significant residual risk.
When an adversary rephrases a known injection to preserve its semantic
intent while altering surface-level features, pattern-matching defenses
fail to recognize the attack. Our evaluation shows that the CIF
firewall's TF-IDF and embedding-based classifiers achieve 89\% detection
on direct injections but only 72\% on semantically equivalent
reformulations. This gap is not unique to CIF; it reflects a fundamental
limitation of feature-based detection that also affects commercial tools
such as Lakera Guard and LLM Guard. Incorporating large language
model-based semantic analysis into the firewall classification pipeline
represents the most promising mitigation, though it introduces
additional latency and cost trade-offs.

Progressive drift attacks exploit the tension between detection
sensitivity and observation window length. An attacker who modifies
beliefs by amounts below the drift detection threshold
(\(\epsilon_{\text{drift}}\)) in each interaction can accumulate
significant deviation over many rounds. Our sliding window detector
catches abrupt changes effectively but misses gradual drift that stays
within per-step bounds. Extending the observation window improves drift
detection but increases response latency proportionally; a 10x window
extension, for example, requires maintaining 10x more historical state.
Adaptive baseline approaches that adjust thresholds based on cumulative
deviation, rather than per-step thresholds alone, offer a promising
direction explored in our supplementary algorithms.

Orchestrator compromise falls outside our current threat model, which
assumes an honest orchestrator. If the orchestrator itself is
compromised, it can selectively disable defenses, suppress alerts, or
manipulate the trust matrix directly. Multi-orchestrator architectures
with cross-verification provide a potential mitigation, effectively
applying Byzantine consensus principles to the orchestration layer
itself. This extension requires careful attention to consistency
guarantees and is an active area of our ongoing research.
\end{block}

\begin{block}{Scaling Beyond Ten Agents}
\protect\phantomsection\label{scaling-beyond-ten-agents}
Our evaluation focused on systems with 3--10 agents, reflecting current
production norms. Three bottlenecks constrain scaling:

\begin{enumerate}
\item
  \textbf{Consensus latency}: Byzantine consensus requires all-to-all
  communication (\(O(n^2)\) messages per round), reaching 4.2s at 100
  agents (\cref{tab:agent-scaling}). Beyond \$\sim\$50 agents,
  hierarchical consensus (partitioning agents into committees with
  inter-committee agreement) or random committee selection is required.
\item
  \textbf{Provenance chain depth}: Deeply nested delegation hierarchies
  slow verification as each link requires cryptographic validation. At
  delegation depth \$\textgreater\$10, provenance pruning (retaining
  only chain endpoints with Merkle commitments for intermediate links)
  becomes necessary.
\item
  \textbf{Memory}: Trust matrix storage (\(O(n^2)\)) and belief history
  (\(O(n \cdot t)\)) dominate. At 100 agents, peak memory reaches 1.6 GB
  (\cref{tab:agent-scaling}). Sliding-window belief history and sparse
  trust representations can mitigate this.
\end{enumerate}
\end{block}

\begin{block}{Generalization Beyond the Evaluated Corpus and
Architectures}
\protect\phantomsection\label{generalization-beyond-the-evaluated-corpus-and-architectures}
Our evaluation, while comprehensive within its scope, faces four
categories of generalization limitations that practitioners should
consider when extrapolating our results to their deployments.

\textbf{Corpus Temporal Bias.} Our attack corpus (950 attacks across
four categories) was assembled from published jailbreak datasets, custom
adversarial prompts, red team exercises, and synthetic generation, but
it necessarily reflects attack techniques known at the time of
evaluation. The adversarial landscape evolves continuously; novel attack
categories---such as attacks exploiting emergent behaviors in
large-scale agent swarms, attacks that manipulate shared environmental
state rather than direct communication channels, or attacks leveraging
model-specific vulnerabilities discovered after our evaluation---may
expose detection gaps not captured by our current corpus. Detection
rates should therefore be interpreted as lower bounds on CIF's
protective capability under the current threat landscape, rather than as
guarantees of future performance against novel adversarial techniques.

\textbf{Scale Testing Limits.} Our evaluation focused on systems with
3--10 agents, which reflects the majority of current production
deployments but does not validate CIF performance at larger scales.
Emerging applications in simulation, scientific discovery, and
autonomous operations may involve 50--100+ agents, where consensus
latency (quadratic in agent count), provenance chain depth, and belief
history storage may introduce both performance degradation and novel
attack surfaces not present at smaller scales. The composition theorems
from Part 1 hold mathematically at any scale, but practical
implementation constraints may force approximations (e.g., hierarchical
consensus, provenance pruning) whose security implications remain
untested.

\textbf{LLM Behavior Variance.} Our simulation-based architecture
adapters model topology and communication patterns but do not execute
actual LLM inference. Language model behavior varies across model
families (GPT-4, Claude, Gemini, Llama), fine-tuning approaches,
temperature settings, and prompt formats in ways that may affect attack
success rates and detection efficacy. CIF detection mechanisms rely on
statistical patterns in agent communication; if production LLM outputs
exhibit different distributional characteristics than our simulation
assumptions, detection thresholds calibrated on our evaluation may
require adjustment. Future work should validate CIF on live inference
with diverse model backends.

\textbf{Architecture Sampling.} The six architectures in our evaluation
(hierarchical orchestrator, peer-to-peer, role-based teams, state
machine, pipeline, and hybrid) represent dominant deployment patterns
but do not exhaust the space of possible multiagent coordination
topologies. Novel architectural paradigms---such as mixture-of-experts
agents, recursive self-improvement loops, or dynamically reconfiguring
topologies---may present vulnerability patterns not captured by our
selected architectures. The composition algebra (Part 1) provides a
principled basis for analyzing new architectures, but empirical
validation on each novel topology is necessary before claiming coverage.

The defense evolution strategy outlined in our adaptive defenses
discussion and the practical \textbf{Risk Assessment Framework} in Part
3 provide concrete strategies for managing these residual generalization
risks through ongoing corpus expansion, periodic defense retraining, and
architecture-specific validation.
\end{block}

\begin{block}{Simulation vs.~Live Deployment Caveats}
\protect\phantomsection\label{simulation-vs.-live-deployment-caveats}
Our simulation-based evaluation approach, while enabling systematic
cross-architecture comparison at scale, introduces important caveats.
The architecture adapters model topology and communication patterns but
do not capture the full complexity of production deployments, including
framework-specific quirks, version-dependent LLM behaviors, or
real-world network conditions. The baseline detection rate of 0.00
reflects the absence of CIF components specifically, not the absence of
all defenses---production frameworks include native safety features
(e.g., Claude Code's permission gating, LangChain's guardrails) that
would provide non-zero baseline protection. Future work should deploy
CIF as middleware on live framework instances to measure marginal
improvement over existing protections.

Additionally, the R\(^2\) values for our scaling regressions (0.994 for
detection time) reflect the controlled simulation environment rather
than the variance typical of production measurements. Practitioners
should expect higher variance in real deployments.
\end{block}

\begin{block}{Observed Cost-Benefit Profile}
\protect\phantomsection\label{observed-cost-benefit-profile}
The ablation data (\cref{tab:component-removal}) reveals a clear
incremental cost-benefit curve. The firewall alone achieves 74\%
detection; adding tripwires raises this to 85\%; the full defense stack
achieves 94\% at 20--25\% latency overhead. The marginal detection gain
per component (\cref{fig:ablation-study}) follows a diminishing-returns
pattern: the firewall contributes \(\Delta\)TPR = +0.13, tripwires
+0.09, provenance +0.07, with subsequent components contributing
\$\leq\$0.06 each.

The Minimal-C configuration (Firewall + Tripwires + Drift Detection)
achieves 90\% detection at 12\% latency overhead
(\cref{tab:minimal-configs}), representing the highest
detection-to-overhead ratio observed. Without trust calculus, lateral
movement attacks in peer-to-peer topologies succeed at rates exceeding
60\% even when other defenses are active. The practical deployment
implications of these findings are explored in Part 3.
\end{block}

\begin{block}{Threats to Validity}
\protect\phantomsection\label{threats-to-validity}
Several threats to validity constrain the generalizability of our
findings. Regarding internal validity, our simulation-based evaluation
models architectural topology and communication patterns but does not
execute actual LLM inference or production framework code. The detection
rates therefore reflect CIF's ability to identify structural attack
patterns rather than its performance against attacks that exploit
specific LLM behaviors or framework vulnerabilities. A follow-up study
deploying CIF as middleware on live framework instances is needed to
establish ecological validity.

External validity is bounded by our selection of six architectures.
While these represent the dominant deployment patterns in current
practice, novel architectural paradigms (such as large-scale swarm
systems or hierarchical mixtures of experts) may present vulnerability
patterns not captured by our evaluation. The 950-attack corpus, though
comprehensive relative to existing benchmarks, cannot represent the full
space of possible cognitive attacks; detection rates should be
interpreted as lower bounds that may decrease when confronting genuinely
novel attack techniques.

Construct validity concerns center on the detection rate metric itself.
A binary detected/undetected classification does not capture partial
detection (e.g., an attack that is flagged but not blocked) or the
severity of successful attacks. Future work should incorporate
severity-weighted metrics and measure time-to-detection alongside binary
classification. Statistical conclusion validity is supported by large
sample sizes, significance testing with Bonferroni correction for
multiple comparisons, and large effect sizes (Cohen's \(d > 0.8\)), but
the controlled simulation environment produces lower variance than
production measurements would exhibit.

Researcher degrees of freedom present a further concern: the framework,
attack corpus, evaluation methodology, and analysis were developed by a
single research group. While we mitigate this through deterministic
reproducibility (fixed seed, public code), pre-registered analysis
protocols (all hypotheses stated before evaluation), and independent
ground-truth labeling (Cohen's \(\kappa = 0.84\) inter-rater agreement),
independent replication by external teams is essential for establishing
the robustness of these findings. We encourage the community to
reproduce our results using the provided scripts and to evaluate CIF
against independently developed attack corpora.
\end{block}
\end{frame}

\begin{frame}{Relationship to Prior Work}
\protect\phantomsection\label{relationship-to-prior-work}
Our empirical results contextualize CIF's contributions relative to the
related work surveyed in \cref{sec:related-work}. Three findings merit
specific comparison.

First, CIF's cognitive firewall achieves 85\% detection on prompt
injection when deployed alone---comparable to published detection rates
for commercial single-agent tools such as NeMo Guardrails
\cite{rebedea2023nemo} and Lakera Guard---but the full CIF stack reaches
96\% by composing the firewall with mechanisms targeting attack vectors
that single-agent tools do not address (trust exploitation, belief
manipulation, coordination). This validates the central thesis that
multiagent security requires defenses beyond input filtering. A
head-to-head comparison on standardized benchmarks (e.g., StrongReject)
remains an important direction for future work.

Second, our trust calculus with \(\delta^d\) decay is, to our knowledge,
the first formally verified bound on delegated trust in LLM-based agent
systems. While classical trust frameworks (FIRE \cite{huynh2006fire},
REGRET \cite{sabater2001regret}) address trust propagation in
distributed systems, none provide the exponential decay guarantee that
our empirical results confirm prevents trust laundering across all six
tested architectures.

Third, CIF's adaptation of Byzantine consensus to semantic content
(beliefs and trust assertions rather than transaction ordering) extends
classical BFT \cite{lamport1982byzantine, castro1999practical} into a
domain where ``Byzantine'\,' behavior manifests as belief poisoning and
coordinated deception. Our 90\% detection rate on coordination attacks
demonstrates practical viability of this adaptation.
\end{frame}

\begin{frame}{Open Research Directions}
\protect\phantomsection\label{open-research-directions}
\begin{block}{Adversarial Retraining and Honeypot Agents}
\protect\phantomsection\label{adversarial-retraining-and-honeypot-agents}
Detection rates inevitably degrade as adversaries learn to evade
deployed defenses---a dynamic well-characterized by the arms race model
in security research and formalized in Part 1's detection degradation
analysis (Section 4). Our current evaluation uses a static attack
corpus, which represents a snapshot of adversarial capability rather
than an evolving threat. Future work should investigate adversarial
retraining of CIF detection mechanisms, where the firewall and anomaly
detectors are periodically retrained on attacks that successfully evade
current configurations. This approach, analogous to adversarial training
in machine learning robustness research, requires careful management of
the retraining loop to prevent catastrophic forgetting of previously
effective detection patterns.

A second promising direction is the deployment of honeypot
agents---intentionally vulnerable agents designed to attract and
characterize novel attack techniques without risking production systems.
By analyzing the attacks directed at honeypot agents, defenders can
identify new attack categories and update detection signatures before
these techniques are deployed against production agents. The formal
framework from Part 1 provides a natural basis for honeypot design: a
honeypot agent can advertise artificially high trust values or
intentionally weak belief validation to attract trust exploitation and
belief manipulation attempts.

Finally, establishing formal safety margins for bounded detection
degradation would allow practitioners to predict the window of
effectiveness for a given defense configuration. If the expected
degradation rate can be bounded, organizations can proactively schedule
defense updates before detection rates fall below acceptable thresholds.
\end{block}

\begin{block}{Collective Invariants for Large-Scale Agent Populations}
\protect\phantomsection\label{collective-invariants-for-large-scale-agent-populations}
As multiagent systems scale beyond the 3--10 agent range evaluated in
this paper, emergent collective behaviors become both a powerful
capability and a significant security concern. Agent collectives may
develop communication patterns, specialization strategies, or
coordination protocols that were not explicitly programmed---some
beneficial, others potentially indicating compromise.

Three concrete research directions address this challenge. First,
\emph{collective invariants} extend CIF's per-agent behavioral
invariants to population-level properties: for example, requiring that
the entropy of the agent interaction graph remains within bounds (sudden
decreases may indicate covert channel formation) or that the
distribution of trust scores across the population does not become
bimodal (which may indicate faction formation by compromised agents).
Second, \emph{emergent behavior fingerprinting} applies graph-theoretic
anomaly detection to the evolving agent interaction network, flagging
topological changes (new cliques, bridge nodes, spectral gap shifts)
that correlate with coordination attacks in our corpus. Third,
\emph{safe emergence boundaries} formalize the conditions under which
emergent behaviors preserve CIF's integrity guarantees---for instance,
proving that if individual agents satisfy trust boundedness (Part 1,
Theorem 4.2), then any emergent coordination pattern among honest agents
also satisfies bounded collective trust, regardless of the specific
strategy that emerges.

These directions connect CIF to the broader complex systems literature
and address the gap between individual-agent security
(well-characterized by our current results) and population-level
security (an open problem as agent counts grow).
\end{block}

\begin{block}{Federated Trust Across Organizational Boundaries}
\protect\phantomsection\label{federated-trust-across-organizational-boundaries}
Current CIF deployment assumes a single operator controlling all agents
within a trust domain. As multiagent systems increasingly span
organizational boundaries---for example, when an enterprise agent
collaborates with agents operated by partners, vendors, or
customers---federated trust management becomes essential.

Federated trust across organizational boundaries requires protocols for
establishing, communicating, and decaying trust between agents that do
not share a common trust authority. The trust calculus from Part 1
provides the mathematical foundation, but practical federation requires
additional mechanisms for trust bootstrapping, cross-domain attestation,
and handling trust domain conflicts. Cross-system provenance
verification presents complementary challenges: verifying the provenance
of information that has passed through agents outside one's trust domain
requires cryptographic techniques (such as verifiable credentials or
zero-knowledge proofs) that go beyond CIF's current provenance tracking.
Regulatory compliance adds a further dimension, as different
jurisdictions impose varying requirements on AI agent autonomy, data
handling, and accountability. A federated CIF deployment must
accommodate these constraints while maintaining coherent security
guarantees across the federation.
\end{block}
\end{frame}

\end{document}
